\documentclass[12pt]{article}
\usepackage{graphicx,amssymb,amsmath,amsthm}
\usepackage{hyperref,url}
\usepackage{verbatim}
\usepackage[labelfont=bf]{caption}
\newcommand{\nocopyright}{
No Copyright\thanks{
The authors hereby waive all copyright
and related or neighboring rights to this work,
and dedicate it to the public domain.
This applies worldwide.
}}
\title{Taut fillings}
\author{Peter Doyle \and Matthew Ellison \and Zili Wang\thanks{
Zili Wang is supported by
the National Natural Science Foundation of China (Grant No. 12301432)
}}
\date{Version A-gpt-01, dated 2026-06-26}
\newtheorem{theorem}{Theorem}
\newtheorem{prop}{Proposition}
\newtheorem{lemma}[theorem]{Lemma}
\newtheorem*{hypo}{Hypothesis}
\newtheorem*{question}{Question}
\newtheorem{corollary}{Corollary}
\newtheorem*{conjecture}{Conjecture}

\def\startproof{{\bf {\medskip}{\noindent}Proof: }}
\def\definition{{\bf {\bigskip}{\noindent}Definition: }}
\def\endproof{$\spadesuit$  \newline}

\newcommand{\proofstart}{{\noindent \bf Proof.\ }}
\newcommand{\note}{{\noindent \bf Note.\ }}
\newcommand{\pfend}{\spadesuit}
\newcommand{\mathproofend}{\quad \qed}
\newcommand{\proofend}{$\quad \qed$}
\newcommand{\goesto}{\rightarrow}
\newcommand{\p}{\;\;\;}
\newcommand{\comma}{\p,}
\newcommand{\pd}{\p.}
\newcommand{\semi}{\p;}
\newcommand{\Rplus}{{\RR_{\geq 0}}}



\newcommand{\beginfigurepos}{\begin{figure}[htbp]}
%
% \fig
% followed by {filename}{label}{caption}
%
\newcommand{\fig}[3]{
\beginfigurepos
\includegraphics[width=370pt]{figures/#1.pdf}
\caption{#3}
\label{#2}
\end{figure}
}
%
% \figsize
% followed by {size}{filename}{label}{caption}
%
\newcommand{\figsize}[4]{
\beginfigurepos
\centerline{
\includegraphics[width=#1]{figures/#2.pdf}
}
\caption{#4}
\label{#3}
\end{figure}
}


\def\ZZ{\mathbf{Z}}
\def\QQ{\mathbf{Q}}
\def\RR{\mathbf{R}}
\def\heartsuit{\mbox{\boldmath{$Z$}}}% 
\def\del{\partial}
\def\dist{\mathrm{dist}}
\def\eps{\epsilon}


\newcommand{\parsec}{\subsection*}
\newcommand{\sect}[1]{\S\ #1}

%\newcommand{\gversus}[2]{\braket{#1\,|\,#2}}
\newcommand{\pair}[1]{\braket{#1}}
\newcommand{\sub}[1]{\braket{#1}}
\newcommand{\nnn}{\mathbf{NSet}}
\newcommand{\zzz}{\mathbf{ZSet}}
\newcommand{\nn}{\mathbf{N}}
\newcommand{\zz}{\mathbf{Z}}

\newcommand{\versus}[2]{\left(#1\,|\,#2\right)}
\newcommand{\gversus}[2]{\left[#1\,|\,#2\right]}
\newcommand{\poly}[1]{\left[#1\right]}

\newcommand{\rising}[1]{\overline{#1}}
\newcommand{\floor}[1]{\left \lfloor #1 \right \rfloor}
\newcommand{\ceiling}[1]{\left \lceil #1 \right \rceil}
\newcommand{\gbinom}{\genfrac{[}{]}{0pt}{}}

\makeatletter
\newcommand{\xRrightarrow}[2][]{\ext@arrow 0359\Rrightarrowfill@{#1}{#2}}
\newcommand{\Rrightarrowfill@}{\arrowfill@\equiv\equiv\Rrightarrow}
\newcommand{\xLleftarrow}[2][]{\ext@arrow 3095\Lleftarrowfill@{#1}{#2}}
\newcommand{\Lleftarrowfill@}{\arrowfill@\Lleftarrow\equiv\equiv}
\newcommand{\xLleftRrightarrow}[2][]{\ext@arrow 3399\LleftRrightarrowfill@{#1}{#2}}
\newcommand{\LleftRrightarrowfill@}{\arrowfill@\Lleftarrow\equiv\Rrightarrow}
\makeatother

\newcommand{\eq}{\Rrightarrow}
\newcommand{\funeq}[1]{\stackrel{#1}\eq}
\newcommand{\funequiv}[1]{\stackrel{#1}\equiv}
\newcommand{\funto}[1]{\stackrel{#1}\to}
\newcommand{\union}{\cup}
\newcommand{\cross}{\times}
\newcommand{\op}{\operatorname}
\newcommand{\kw}{\mathrm}
\newcommand{\escape}{\op{escape}}
\newcommand{\cancel}{\op{cancel}}
\newcommand{\nestuntil}{\operatorname{nestuntil}}
\newcommand{\size}{\operatorname{size}}
\newcommand{\id}{\operatorname{id}}
\newcommand{\create}{\operatorname{create}}
\newcommand{\destroy}{\operatorname{destroy}}
\newcommand{\then}{\triangleleft}
\newcommand{\lp}{\kw{LP}}
\newcommand{\LPvol}{\kw{LPvol}}
\newcommand{\vol}{\kw{vol}}
\newcommand{\flipdist}{\kw{flipdist}}
\newcommand{\tetbar}{\kw{homvol}}
\newcommand{\vsa}{\kw{vsa}}
\newcommand{\cone}{\kw{cone}}
\newcommand{\link}{\kw{link}}
\newcommand{\maxdeg}{\kw{maxdeg}}
\newcommand{\mindeg}{\kw{mindeg}}
\newcommand{\vertices}{\kw{Vert}}
\newcommand{\carrier}{\kw{carrier}}
\newcommand{\nbhd}{\kw{nbhd}}
\newcommand{\adj}{\kw{adj}}
\newcommand{\taut}{\kw{taut}}
\newcommand{\net}{\kw{net}}

\newcommand{\Zvol}{\kw{Zvol}}
\newcommand{\Qvol}{\kw{Qvol}}
\newcommand{\Rvol}{\kw{Rvol}}
\newcommand{\tetvol}{\kw{tetvol}}
\newcommand{\hash}{\#}
\newcommand{\ms}{[}
\newcommand{\msend}{]}

\newcommand{\kwstyle}{\bf}
\newcommand{\IF}{{\ \kwstyle{if}\ }}
\newcommand{\ELSE}{{\ \kwstyle{else}\ }}
\newcommand{\WHILE}{{\kwstyle{while}\ }}
\newcommand{\DO}{{\ \kwstyle{do}\ }}
\newcommand{\RETURN}{{\kwstyle{return}\ }}
\newcommand{\direct}{\oplus}
\newcommand{\pls}{\mathord{+}}
\newcommand{\mns}{\mathord{-}}



\begin{document}

\maketitle

\begin{abstract}
Sleator, Tarjan, and Thurston asked: given a triangulation $\sigma$ of the
$2$-sphere, what is the minimum number of tetrahedra needed to extend
$\sigma$ to a triangulation of the ball?  Call this minimum $\tetvol(\sigma)$.
Let $X(\sigma)$ be either associated integral $2$-cycle, and let
$\Zvol(\sigma)$ be the minimum $L_1$-norm of an integral $3$-chain $M$ with
$\del M=X(\sigma)$.  Certainly $\Zvol(\sigma)\leq \tetvol(\sigma)$.  We prove
that equality holds.  More precisely, any taut filling of $X(\sigma)$ is a
simplicial chain whose support complex is clean and triangulates $B^3$; the
resulting triangulation is freely shellable and flag.  The key input is the
general fact that taut fillings of integral $n$-cycles split under disjoint
union, connected sum, and more generally under almost disjoint union, where
the summands are supported on vertex sets meeting in at most $n+1$ vertices.
\end{abstract}

\section{Overview}
Let $\Delta = \Delta(\Omega)$ be the $|\Omega|\mns 1$-simplex
with vertices $\Omega$,
viewed as the abstract simplicial complex consisting of all subsets of
$\Omega$.
Let $C_n=C_n(\Omega)=C_n(\Delta(\Omega),\ZZ)$
and $Z_n=Z_n(\Omega)=Z_n(\Delta(\Omega),\ZZ)$
be its integral $n$-chains and $n$-cycles.
A \emph{filling} of an $n$-cycle $X \in Z_n$ is any
$n\pls 1$-chain $M \in C_{n+1}$ with $\del M = X$.
(Here and throughout we'll assume $n\geq 1$.)
Let $\Zvol(X)$ be the minimum $L_1$-norm of a filling of $X$.
Call $M \in C_{n+1}$ \emph{taut} if it is an optimal filling
of its boundary, i.e., if $|M|=\Zvol(\del M)$.

For $X \in C_n$ let
$\vertices(X)$ be the set of all the vertices of all the $n$-simplices
to which $X$ assigns non-$0$ weight.
For any taut filling $M$ of $X$ we have $\vertices(M) = \vertices(X)$.
(Cf.\ Proposition \ref{nointernal} below.)
This is why we write $C_n$ and $Z_n$ without reference to $\Omega$.

Call an $n$-cycle $X+Y \in Z_n$ an \emph{almost disjoint union} if
\[
|\vertices(X) \cap \vertices(Y)| \leq n+1
.
\]
This notion generalizes both disjoint union and connected sum
along an $n$-simplex.
Ellison \cite{ellison:gap}
showed the $\Zvol$ adds under almost disjoint union:
$\Zvol(X+Y) = \Zvol(X) + \Zvol(Y)$.
This means that \emph{some} taut filling $M$ of $X$ splits into a sum
$M=M_X+M_Y$ of taut fillings of $X,Y$.
Here we amplify Ellison's Corollary 2 to show (Theorem \ref{th1})
that when $n \geq 2$, \emph{any} taut filling of $X+Y$ splits.

In place of integral chains we can use chains with
coefficients in $\QQ$ or $\RR$.
Evidently $\Rvol \leq \Qvol$.
But computing $\Rvol$ is
a linear programming problem with rational coefficients,
so in fact $\Qvol = \Rvol$.
Ellison used LP duality to prove that
$\Qvol$ adds under almost
disjoint union.
Now for $n \geq 2$ we get the stronger result
(Corollary \ref{cor1})
that taut $\QQ$-fillings split:
Multiplying a taut
$\QQ$-filling $M$ of $X$ by a common denominator $q$
yields a taut $\ZZ$-filling $qM$ of $qX$:
Split the $\ZZ$-filling $qM$ and divide by $q$ to get a splitting
of the $\QQ$-filling $M$.

Now let $\sigma$ be a simplicial triangulation of $S^2$;
let $X(\sigma) \in Z_2(\Delta(\vertices(\sigma)),\ZZ)$
be either one of the 2-cycles
that arises from $\sigma$ by orienting its $2$-simplices;
and write $\Zvol(\sigma) = \Zvol(X(\sigma))$.
Let $\tetvol(\sigma)$ be the minimum number of $3$-simplices
required to extend $\sigma$ to a simplicial triangulation $\tau$ of $B^3$.
We will show that
\[
\Zvol(\sigma) = \tetvol(\sigma)
.
\]
Certainly $\Zvol \leq \tetvol$,
because any extension $\tau$ induces a filling
of $X(\sigma)$.
Theorem \ref{th2} states the stronger fact that every taut filling is clean and has support complex a triangulated ball.  The proof of Theorem \ref{th2} relies on Theorem \ref{th1}.

We don't know what happens for spheres of dimension 3 or greater.

\section{Background}

Our interest in $\Zvol$ stems from the work
of Sleator, Tarjan, and Thurston
\cite{stt:jams}.
They observed that coning from a vertex of maximal degree
shows that a $v$-vertex
triangulation $\sigma$
of $S^2$ has $\tetvol \leq 2v-4-\maxdeg(\sigma)$.
If $v\geq13$ $\maxdeg(\sigma) \geq 6$, so $\tetvol \leq 2v-10$.
They
produced examples for which $\tetvol = 2v-10$, provided
$v$ is larger than some unspecified bound,
and conjectured that such examples should exist for any $v \geq 13$.
Their approach was to 
realize $\sigma$ as an ideal hyperbolic polyhedron
with volume $2v V_0-O(\log(v))$,
where $V_0$ is the volume of an
equilateral ideal tetrahedron,
which is maximal.
This implies $\tetvol \geq 2v-O(\log(v))$.
From there they bootstrapped their way up to showing
that $\tetvol=2v-10$.
They suggested
\cite[p. 697]{stt:jams}
that in fact $\Qvol= 2v-10$,
which would immediately imply $\tetvol=2v-10$.
Mathieu and Thurston \cite{mt}
produced a class of examples,
different from the examples of STT,
for which they could show that $\Qvol = 2v-10$,
still under the assumption that $v$
is sufficiently large.
In \cite{dew:filling} we produced examples for any $v \geq 13$ with
$\Qvol=2v-10$, confirming the conjecture of Sleator, Tarjan, and Thurston.
Theorem \ref{th2} here
tells us that whenever we have $\Qvol=2v-10$,
or just $\Zvol=2v-10$, any taut filling
must arise from a triangulation of the ball.

About $\Qvol$ versus $\Zvol$.
In \cite{dew:filling}
we described triangulations of $S^2$ with $\Qvol < \Zvol$.
In all such examples that we know of
with $\maxdeg \leq 6$, the gap between $\Qvol$ and $\Zvol$ is $<1$,
so that
$\lceil \Qvol \rceil = \Zvol$.
Since $\Qvol$ and $\Zvol$
add under connected sum,
the gap between them can be arbitrarily
large
(Ellison \cite{ellison:gap}),
but taking the connected sum pushes $\maxdeg$ above $6$.



\section{Taut fillings}

We can think of an $n$-chain $X \in C_n$ as a multiset of
non-cancelling oriented simplices:
\[
X = \sum_{t \in X} t
.
\]
In this sum $t$ denotes an oriented simplex
\[
t=[x_0,\ldots,x_n]
= [x_{\pi(0)},\ldots,x_{\pi(n)}]
,\;\;\mbox{$\pi$ an even permutation}
,
\]
and each unoriented simplex contributes a number of terms corresponding
to its multiplicity.

The size of $M$ as a multiset is its $L_1$-norm $|M|$.
Write $U \subset M$
if $U$ is a sub-multiset of $M$.
This happens just when
$|M| = |U| + |M-U|$.

\begin{prop} \label{subtaut}
If $M$ is taut and $U \subset M$ then $U$ is taut.
\end{prop}

\proofstart
This is clear, but let's drag it out.
If $U$ is not taut, there is a smaller filling $V$ of $\del U$.
Take $M^\prime = M-U+V$.
In terms of multisets, this means that we substitute $V$ for $U$,
and then do any required cancellation of oppositely oriented simplices
of $M-U$ and $V$.
We have
\[
\del M^\prime = \del M - \del U + \del V = \del M
\]
and
\[
|M^\prime| \leq |M-U| + |V| = |M| - |U| + |V| < |M|
,
\]
contradiction.
\proofend

\section{Coning}
For $x \in \Omega$, $U \in C_n$ let
\[
\nbhd(x,U) = \sum_{ t \in U: x \in \vertices(t) } t
,
\]
\[
\deg(x,U) = |\nbhd(x,U)|
,
\]
\[
\maxdeg(U) = \max_x \deg(x,U)
.
\]
The \emph{cone from $x$ to $U$}
is the $(n+1)$-chain
\[
\cone(x,U) = 
\sum_{t \in U} \adj(x,t)
,
\]
consisting of all the non-trivial oriented
$(n+1)$-simplices $\adj(x,t)=[x x_0\ldots x_n]$ obtained
by adjoining $x$ to  $t = [x_0\ldots x_n] \in U$.
If $x \in \vertices(t)$ then $t$ doesn't contribute to the sum,
so $|\cone(x,U)| = |U| - \deg(x,U)$.

If $U$ is closed then
$\del\, \cone(x,U) = U$, so

\begin{prop} \label{maxdeg}
For any $X \in Z_n$
\[
\Zvol(X) \leq |X| - \maxdeg(X)
\mathproofend
\]
\end{prop}

If $U \in Z_n$ and $x \notin \vertices(U)$ then
$\deg(x,U)=0$ and
$|\cone(x,U)| = |U|$.
In this case we call $\cone(x,U)$ a \emph{complete cone}.
By Proposition \ref{maxdeg} a non-trivial complete cone is not taut.
(This is a long-winded way of saying that instead of coning from
$x \notin U$, we could have coned from some $x \in U$.)
Thus:
\begin{prop}\label{nocone}
If $M$ is taut it contains no non-trivial complete cone.
\proofend
\end{prop}

Call $x \in \vertices(M) \setminus \vertices(\del M)$ an
\emph{internal vertex} of $M$.
If $x$ is internal to $M$ then $\nbhd(x,M)$ is a complete cone, so:
\begin{prop} \label{nointernal}
If $M$ is taut then
it has no internal vertices.
\proofend
\end{prop}

\section{Almost disjoint unions}

Recall that if $X,Y \in Z_n$ we call $X+Y$ an almost disjoint union
if
\[
|\vertices(X) \cap \vertices(Y)| \leq n+1
.\]
The most interesting special case is a connected sum,
where $t=[c_0,c_1,\ldots,c_n]$ occurs once in $X$ and
$-t=[c_1,c_0,\ldots,c_n]$ occurs once in $Y$.
For example, if $n=1$ with $X$ a cycle of length $p$
and $Y$ a cycle of length $q$
the connected sum $X+Y$ is
a cycle of length $p+q-2$.
Here
$\Zvol$ adds, because
$\Zvol(X)=p-2$,
$\Zvol(Y) = q-2$,
and
\[
\Zvol(X+Y) = (p+q-2)-2 = \Zvol(X)+\Zvol(Y)
.
\]
But not every filling of $X+Y$ splits, because a taut filling
may contain any 2-simplex with vertices in $\vertices(X+Y)$.

We want to show that fillings do split when $n \geq 2$.

For $A \subset \Omega$, $p \in A$ define
\[
\pi_{A,p}: \Omega \to A
\]
\[
\pi_{A,p}(x) = x \IF x \in A \ELSE p
\]
and let
\[
K_*(A,p): C_*(\Omega) \to C_*(A)
\]
be the induced chain map.
This is a projection of $C_*(\Omega)$ onto $C_*(A)$.

Let $A,B$ be finite vertex sets sharing the vertices
$C = A \cap B$,
and let $c=|C|$.
Observe that
$C_c(C)$ and $Z_{c-1}(C)$ are trivial,
because $\Delta(C)$ contains no $c$-simplex,
and (up to orientation) only one $c\mns 1$-simplex,
with nothing to cancel its boundary.

Let
\[
C_*(A,B) = C_*(A) \direct C_*(B)
.
\]
For $p,q \in C$ define the mapping
\[
g_*(p,q) = K_*(A,p) \direct K_*(B,q):
C_*(A \cup B) \to C_*(A,B)
.
\]

\begin{prop} \label{recover}
Let $(X,Y) \in C_n(A,B)$.
If $|A \cap B| \leq n$
we can recover $(X,Y)$ from $X+Y$.
If $(X,Y) \in Z_n(A,B)$ this holds also when $|A \cap B|  = n+1$.
\end{prop}

\proofstart
Let $C=A \cap B$.
We can assume $|C| \geq 1$.
(Add a brand new point to $C$ if necessary.)
We claim that
for any $p,q \in C$
(not necessarily distinct) we have
\[
g_n(p,q)(X+Y)=(X,Y)
.
\]
Indeed,
\[
K_n(A,p)(X+Y) = X + K_n(A,p)(Y)
.
\]
The second term belongs to $C_n(C)$,
which is trivial when $|C| \leq n$.
If $Y \in Z_n(B)$ the second term belongs to $Z_n(C)$,
which is trivial when $|C| \leq n+1$.
\proofend


\begin{theorem} \label{th1}
If $|A \cap B| \leq n+1$,
for all $(X,Y) \in Z_n(A,B)$ we have
\[
\Zvol(X+Y) = \Zvol(X)+\Zvol(Y)
.
\]
And as long as $n \geq 2$,
for any $M \in \taut(X+Y)$ 
we have $M = M_X + M_Y$
with $M_X \in \taut(X)$, $M_Y \in \taut(Y)$.
\end{theorem}

\note
Ellison \cite{ellison:gap}
proved the additivity of $\Zvol$.
What's new here is the splitting of taut fillings when $n\geq 2$.
This result resembles Theorem 6.2 of
Pournin and Wang
\cite{pw:flip}
about flip paths.
Like theirs, our proof uses a variation on the normalization technique
of STT
\cite[Lemma 7]{stt:jams}.
It would be nice to fit these results under one roof.
\medskip

\proofstart
Again let $CV = A \cap B$.
We can assume $|C| = n+1$,
as this is the hardest case.
And we might as well go ahead and take $n=2$, $|C|=3$,
as this case illustrates all the issues.

Take any $M \in \taut(X+Y) \subset C_3(A \union B)$.
Pick distinct points $p,q \in C$,
and let
\[
(M_X,M_Y) = g_{3}(p,q)(M)
.
\]
(For now we'll suppress the dependence of $M_X,M_Y$ on $p,q$.)

Because $g_*(p,q)$ is a chain map we have
$\del_{3} M_X = X$, $\del_{3} M_Y = Y$:
\[
(\del_{3} M_X, \del_{3} M_Y)
=
\del_{3} (g_{3}(p,q)(M))
=
g_2(p, q)(\del_{3} M)
=
g_2(p,q)(X+Y)
=
(X,Y)
.
\]

We want to show that under the map $g_3(p,q)$ any $t \in M$
dies either in $M_X$ or $M_Y$,
because then we'll get additivity of $\Zvol$ from
\[
\Zvol(X+Y) = |M|
\geq |M_X| + |M_Y|
\geq \Zvol(X) + \Zvol(Y)
\geq \Zvol(X+Y)
.
\]

Let's call a $3$-simplex a `tet' for short.
Say that a tet $t \in M$ has type $CCXY$ if $t= \pm[c_1c_2x_1y_1]$
for $c_1,c_2 \in C$, $x_1 \in A \setminus C$, $y_1 \in B \setminus C$.
Similarly for types $XXXX$, $CXXX$, $CXXY$, $XXXY$, etc.
The first two are pure $X$ cases, meaning that they live in $C_3(A)$;
the last two are hybrid cases.

Any $XX..$ tet dies in $M_Y$;
any $YY..$ tet dies in $M_X$.

The remaining cases to check are
the pure cases $CCCX$, $CCCY$,
and the hybrid case $CCXY$.
A $CCCX$ tet $t=\pm[c_1c_2c_3x_1]$ dies in $M_Y$ because
$K_3(B,q)(t)=[c_1c_2c_3q]$,
which vanishes because $q \in C = \{c_1,c_2,c_3\}$ produces a repeated vertex.
Likewise a $CCCY$ tet dies in $M_X$.
The more interesting case is $CCXY$:
The key is that since $|C|=3$, $\{c_1,c_2\}$ cannot be disjoint from $\{p,q\}$,
so it dies on one side or the other.

So 
\[
|M| = |M_X(p,q)| + |M_Y(p,q)|,
\]
and $\Zvol$ adds.
Note how we're now emphasizing the possible dependence of $M_X,M_Y$ on $p,q$.
When $n=1$ the choice of $p,q$ can indeed make a difference:
We can get a different pair $(M_X,M_Y)$ if we switch
$p$ and $q$.
(Think about the connected sum of two cycle graphs.)

But when $n \geq 2$ we will now show that all tets are pure $X$ or pure $Y$,
so $M_X$ consists of all the pure $X$ tets, and ditto for $M_Y$.
This will make
$M = M_X + M_Y$ with $M_X \in \taut(X)$, $M_Y \in \taut(Y)$.

Again our test case is $n=2$.

The key observation is that, now that we know that every tet dies on
one side or the other, we know that no tet can die on both sides,
because that would make $|M| > \Zvol(X) + \Zvol(Y)$.

An $XXYY$ tet would die on both sides, so there can be none of these.

Any $pqXY$, $pXXY$, or $qXYY$ would die on both sides, and $p,q$ are arbitrary,
so this rules out all $CCXY,CXXY,CXYY$.

The only remaining hybrids are $XXXY$ and $XYYY$.
Let's rule out $XXXY$, leaving $XYYY$ to symmetry.

Fix any $y_0 \in B \setminus C$, and suppose there is
some $XXXy_0$ tet in $M$.
Let $U \subset M$ consist of all such $XXXy_0$ tets.
We claim that $y_0 \notin \vertices(\del U)$.
For let $s$ be any 2-simplex of the form $XXy_0$,
the only kind of 2-simplex containing $y_0$ that could belong
to $\del U$.
Any $t \in M$ of which $s$ or $-s$ is a face must be a hybrid tet,
and the only possibility is $XXXy_0$, so $t \in U$.
Because $s$ has no support in $X+Y=\del M$, the signed multiplicity of $s$
vanishes in $\del M$,
hence also in $\del U$,
because the tets of $M$ with $\pm s$ as a face all belong to $U$.
So $y_0 \notin \vertices(\del U)$,
making $U$ a complete cone on $y_0$,
contradicting \ref{nocone}.
So there is no such $XXXy_0$ tet, ruling out $XXXY$,
and with it $XYYY$.

So there are no hybrid tets, meaning that $M$ splits, as claimed.
\proofend

The splitting of taut fillings depends on the assumption $n\geq 2$.
Let's consider what goes wrong when $n=1$.
Here the hybrid types are $CXY$, $XXY$, $XYY$.
Neither $pXY$ nor $qXY$ dies on both sides, so we can't rule out $CXY$.
Failing that, the $XXy_0$ tets needn't form a complete cone,
because $[x_0x_1y_0]$ can continue across $[x_1 y_0]$ to $-[p x_1 y_0]$,
so we can't rule out $XXY$ either.

The foregoing proof works equally well over $\QQ$, providing we
permit ourselves to work with fractional multisets.
Alternatively, we can clear denominators as in the introduction
above.
Either way, we have:

\begin{corollary} \label{cor1}
$\Qvol$ adds under almost disjoint union,
and for $n \geq 2$
taut $\QQ$-fillings split.
\proofend
\end{corollary}

\section{Triangulations}

We turn now to filling simplicial triangulations of $S^2$.
We first fix the combinatorial language.

An \emph{$n$-simplex} is a set of size $n+1$.
Its faces are its subsets, regarded as simplices of dimensions
$-1,0,\ldots,n$, where the empty simplex has dimension $-1$.
A \emph{simplicial complex} $K$ is a finite collection of simplices closed
under taking faces: if $t\subset s\in K$, then $t\in K$.

For a simplex $s\in K$, define the \emph{link}
\[
\link(s,K)=\{t: t\cap s=\emptyset,\ s\cup t\in K\}.
\]
This is again a simplicial complex.  Except when $s=\emptyset$, it is not
literally a subcomplex of $K$: it records the faces complementary to $s$ in
simplices containing $s$.

We will call an $n$-dimensional simplicial complex \emph{clean} if it is a
normal pseudomanifold with boundary, in the following concrete sense.
First, it is pure: every simplex belongs to some $n$-simplex.  Second, every
$(n-1)$-simplex lies in either one or two $n$-simplices; those lying in only
one form the boundary.  Third, for every simplex $s$ of dimension
$0,\ldots,n-2$, the link $\link(s,K)$ is connected.  Thus in dimension $3$
we are requiring the usual triangle-count condition together with connected
edge links and connected vertex links.  We do not impose connectedness of
$K$ itself as a separate convention; in the applications below the complexes
under discussion are connected.

The boundary of a clean complex is clean and closed, and links in a clean
complex are clean of the corresponding dimension.

If $N$ is a manifold, a simplicial complex $K$ is a \emph{simplicial
triangulation} of $N$ if its geometric carrier is homeomorphic to $N$.  When
$N$ is oriented, a choice of orientation of the triangulation gives an
integral cycle $X(K)$; in the case $N=S^2$ the two choices differ by sign, and
nothing below depends on which one is chosen.

This notion is the ordinary simplicial one.  We do not allow two distinct
simplices with the same vertex set, nor do we allow self-identifications of a
simplex along its own faces.

We continue to regard an $n$-chain $M\in C_n$ as a multiset of oriented
$n$-simplices.  We call $M$ \emph{simplicial} if no unoriented simplex occurs
with multiplicity greater than one.  If $M$ is simplicial, let $K(M)$ denote
the pure simplicial complex generated by the unoriented simplices occurring
in $M$.  A filling $M$ is a chain, not a complex; we will nevertheless call
$M$ \emph{clean} if $M$ is simplicial and $K(M)$ is a clean complex.

\section{Filling a triangulation of the 2-sphere}

\begin{theorem} \label{th2}
Let $\sigma$ be a simplicial triangulation of $S^2$, and let $M$ be a taut
filling of $X(\sigma)$.  Then $M$ is clean, and $K(M)$ is a simplicial
triangulation of $B^3$.
\end{theorem}

\proofstart
Call $(\sigma,M)$ a \emph{filling pair} if $\sigma$ is a simplicial
triangulation of $S^2$ and $M$ is a taut filling of $X(\sigma)$.
Call such a pair \emph{good} if $M$ is clean and $K(M)$ is a simplicial
triangulation of $B^3$.  We prove that every filling pair is good.

Suppose not, and choose a bad filling pair $(\sigma,M)$ for which $|M|$ is
minimal.  Let $v,e,f$ be the numbers of vertices, edges, and faces of
$\sigma$; then $e=3v-6$ and $f=2v-4$.  The tetrahedron boundary is the base
case, so $v>4$.  Since taut fillings split under connected sum, we may assume
that $\sigma$ is prime.  In particular, the only consequence we shall use is
that $\sigma$ has no vertex of degree $3$.

A tet $t\in M$ is \emph{eligible} if two of its oriented faces occur as faces
of $\sigma$ with the boundary orientation.  It cannot have more than two such
faces, because that would give a degree-$3$ vertex of $\sigma$.  By
Proposition \ref{maxdeg},
\[
 |M|=\Zvol(\sigma)\leq f-\maxdeg(\sigma).
\]
For each support tet $u$ of $M$, let $A(u)$ be the set of faces of $\sigma$
which occur as properly oriented faces of $u$.  These sets cover the faces of
$\sigma$, since each boundary face has nonzero coefficient in $\del M$ with
the sign specified by $X(\sigma)$.  Moreover $|A(u)|\leq 2$.  Reveal the
support tets one at a time, and count how many previously uncovered boundary
faces each reveals.  If $r$ tets reveal two new boundary faces, then
\[
 f\leq |M|+r.
\]
Hence $r\geq \maxdeg(\sigma)$.  These $r$ tets are eligible, and the pairs of
new boundary faces assigned to them are disjoint.  We shall use only the
existence of two such eligible tets, and the fact that for any fixed boundary
face there is an eligible tet not containing it.

Choose an eligible tet $t$ with boundary faces
\[
 s_1=[a,b,c],\qquad s_2=[b,a,d].
\]
Removing $t$ replaces these two boundary faces by
\[
 [c,d,b],\qquad [d,c,a],
\]
which is the boundary flip replacing the edge $ab$ by the edge $cd$.
Let the resulting boundary complex be denoted by $\sigma_t$.  There are two
possibilities.

In the first case, $cd$ was not an edge of $\sigma$, and $\sigma_t$ is again
a simplicial triangulation of $S^2$.  Then $M-t$ is a taut filling of
$\sigma_t$, so by minimality $(\sigma_t,M-t)$ is good.  Thus $M-t$ is clean
and $K(M-t)$ triangulates $B^3$.  Gluing $t$ back along the two exposed
boundary faces gives $K(M)$, provided $t$ was not already present in $M-t$.
If $t$ were already present, then $M$ would contain the same unoriented tet
twice.  Choose another eligible tet $u$ whose two boundary faces are disjoint
from those of $t$.  Removing $u$ leaves a smaller filling in which the doubled
tet $t$ is still present, so the smaller pair is still bad, contradicting the
minimality of $(\sigma,M)$.  Hence no such doubling occurs, and in this case
$M$ is good.

In the second case, $cd$ was already an edge of $\sigma$.  Then $\sigma_t$ is
an almost disjoint union of two triangulated spheres, say $\sigma_1$ and
$\sigma_2$, joined along the edge $cd$.  The filling $M-t$ need not be clean
as one complex: the link of the common edge is disconnected.  But by
Theorem \ref{th1}, the taut filling $M-t$ splits as
\[
 M-t=M_1+M_2,
\]
where $M_i$ is a taut filling of $\sigma_i$.  By minimality, each
$(\sigma_i,M_i)$ is good.  Thus $K(M_1)$ and $K(M_2)$ are clean triangulated
balls.  Gluing these two balls together with the bridging tet $t$ gives
$K(M)$, again unless $t$ was already present in $M-t$.  The same second
eligible tet argument rules this obstruction out.  Therefore $M$ is good in
this case as well.

It remains only to spell out why this gluing preserves cleanness.  Away from
the two faces along which $t$ is attached, the links and face incidences are
unchanged from the smaller clean ball or balls.  Along the gluing faces, each
boundary triangle becomes incident to exactly two tets.  The new local edge
and vertex links are obtained from connected intervals or disks by attaching
one simplex along a nonempty connected part, so they remain connected.  The
only possible failures would be a repeated tet, a triangle used by more than
two tets, or a disconnected edge or vertex link already visible after removing
a disjoint eligible tet.  Such a persistent failure would make the smaller
pair bad, contradicting the chosen minimality.  Hence the reassembled complex
is clean and triangulates $B^3$.

This contradiction proves the theorem.
\proofend

Thanks to this theorem, we may take the liberty of saying that a taut filling
of $\sigma$ is a simplicial triangulation $\tau$ of $B^3$.

\section{Shelling}

We next show that a taut filling may be shelled from any chosen initial tet.
A \emph{shelling} of a simplicial triangulation $\tau$ of $B^3$ is an ordering
$(t_1,\ldots,t_{|\tau|})$ of its tets such that, for every $k$, the subcomplex
generated by $t_1,\ldots,t_k$ is again a simplicial triangulation of $B^3$.
For $k>1$, the new tet $t_k$ may meet the previous complex in one, two, or
three faces.  In our situation the three-face case is excluded by the absence
of internal vertices.  We call $\tau$ \emph{freely shellable} if every tet of
$\tau$ can be chosen as the first tet of some shelling.

\begin{theorem} \label{th3}
If $\sigma$ is a simplicial triangulation of $S^2$, any taut filling $\tau$ of
$\sigma$ is a freely shellable simplicial triangulation of $B^3$.
\end{theorem}

\proofstart
This follows by turning the induction in Theorem \ref{th2} into a recursive
``shucking'' procedure.  Fix the tet $t_1$ which is to be first in the final
shelling; equivalently, we will remove all other tets first and remove $t_1$
last.

If $\tau$ consists only of $t_1$, there is nothing to prove.  If a degree-$3$
vertex belongs to a tet $t\neq t_1$, remove that tet and shuck the remaining
filling, still ending with $t_1$.  If the only available degree-$3$ vertex
belongs to $t_1$, let $t_2$ be the tet of $\tau-t_1$ to which $t_1$ is
attached; shuck $\tau-t_1$ ending with $t_2$, and then put $t_1$ first when
reading the removals backwards.

If there is no degree-$3$ vertex, the eligible-tet count supplies at least
$\maxdeg(\sigma)\geq 4$ disjoint eligible tets.  Choose one sharing no face
with $t_1$.  If removing it leaves one sphere boundary, recurse on the smaller
filling.  If removing it splits the boundary into two conjoined spheres, split
the filling as in Theorem \ref{th1}, shuck the side not containing $t_1$ first,
then remove the bridging tet, and finally shuck the side containing $t_1$.
Reading this removal order backwards gives a shelling beginning with $t_1$.
\proofend

\section{Flag complex}

A simplicial complex is a \emph{flag complex} if every clique in its
$1$-skeleton occurs as a simplex.  A triangulation of $S^2$ is flag exactly
when it is prime and is not the boundary of a tetrahedron.

We can decompose any triangulated $2$-sphere as a connected sum of prime
components.  Taut fillings split along this connected-sum decomposition, so
it is enough to understand the prime pieces; nevertheless the conclusion
below is stated for arbitrary $\sigma$.

\begin{theorem} \label{th4}
If $\sigma$ is a simplicial triangulation of $S^2$, any taut filling $\tau$ of
$\sigma$ is a flag complex.
\end{theorem}

\proofstart
We must rule out three taboo configurations in $\tau$:
\begin{enumerate}
\item the three edges of a triangle $\{A,B,C\}$ without the face $ABC$;
\item the six edges of a $K_4$ on $\{A,B,C,D\}$ without the tet $ABCD$;
\item the ten edges of a $K_5$ on $\{A,B,C,D,E\}$.
\end{enumerate}
If the first two configurations are ruled out, the third is ruled out as
well: a $K_5$ would force all five of its $K_4$ subcomplexes to occur as tets,
forming the boundary of a $4$-simplex, which cannot occur inside a triangulated
$3$-ball.

Suppose there is a minimal counterexample.  If $\sigma$ has a vertex of
degree $3$, then $\sigma$ is a connected sum with one summand the boundary of
a tetrahedron.  Removing the corresponding tet from $\tau$ gives a smaller
counterexample, a contradiction.  Thus we may assume that $\sigma$ has no
vertex of degree $3$.

Then the eligible-tet count gives at least $\maxdeg(\sigma)$ disjoint eligible
tets.  Removing one performs a boundary edge flip.  The new boundary either
remains a triangulated sphere or becomes an almost disjoint union of two
spheres joined along an edge.  In the first case, preserving the taboo
configuration gives an immediate smaller counterexample.  In the second case,
a taboo $K_3$ or $K_4$ in the almost disjoint union lies entirely on one side
or the other, again giving a smaller counterexample after splitting the taut
filling.

For a missing triangle, only one edge disappears when we remove an eligible
tet: the boundary edge being flipped.  There are at least four eligible flips,
and only the three edges of the taboo triangle need to be avoided.  Thus some
flip preserves the missing triangle, contradicting minimality.

For a missing tet on vertices $A,B,C,D$, there are six edges to avoid.  Since
missing triangles have already been ruled out, all four faces
$ABC,ABD,ACD,BCD$ occur.  At most two of these faces can lie in the boundary
$\sigma$, because otherwise $\sigma$ would have a degree-$3$ vertex.  Treating
the octahedral boundary as the small exceptional case, we may assume
$\maxdeg(\sigma)\geq 5$.  Hence there are at least five disjoint eligible
flips.  If five of the six $K_4$ edges are not simultaneously boundary edges,
one of these flips avoids the taboo $K_4$.  In the remaining case, say the
only interior edge is $AD$, so the boundary contains the faces $ABC$ and
$BCD$.  At most two of the disjoint eligible tets can use one of these two
faces, leaving another eligible flip which preserves the $K_4$.  Again we get
a smaller counterexample, contradiction.
\proofend

\begin{thebibliography}{5}

\bibitem{dew:filling}
Peter Doyle, Matthew Ellison, and Zili Wang,
\emph{Filling a triangulation of the 2-sphere}, 2024.

\bibitem{ellison:gap}
Matthew Ellison,
\emph{Lower bounds and integrality gaps in simplicial decomposition}, 2024.

\bibitem{mt}
Claire Mathieu and William Thurston,
\emph{Rotation distance using flows}, 1992.

\bibitem{pw:flip}
Lionel Pournin and Zili Wang,
\emph{Strong convexity in flip-graphs}, 2021.

\bibitem{stt:jams}
Daniel D. Sleator, Robert E. Tarjan, and William Thurston,
\emph{Rotation distance, triangulations, and hyperbolic geometry},
J. Amer. Math. Soc. \textbf{1} (1988), 647--681.

\end{thebibliography}
\end{document}
